% !TeX root = ../main.tex
\section{Note transcription}
\label{sec:notes}

\subsection{What the stage must do}

Input: the guitar stem (or the mix). Output: a list of notes, each with onset,
offset, pitch and a loudness estimate. This is \term{automatic music
transcription} (AMT) restricted to one instrument.

\subsection{Frame-and-onset models}

\begin{keyidea}[: predict onsets and sustain separately]
Onsets and Frames~\cite{hawthorne2018onsets} trains one network to output two
posteriorgrams: where notes \emph{start} and where notes \emph{sound}. A note is
created only at a detected onset and lasts while the frame output stays high.
Separating the two makes repeated notes and sustained notes both readable.
\end{keyidea}

Kong et al.~\cite{kong2021highres} improved this for piano by \emph{regressing}
the exact onset and offset time inside each 10\ms{} frame instead of a yes/no
target, which gives sub-frame timing and tolerates labels that are a little off.
It reached 96.72\% onset F1 on the MAESTRO piano dataset \reported. This
architecture (a CRNN with onset, offset, frame and velocity outputs at 100
frames per second) is the basis of today's best guitar models: Riley et al.\
first adapted it to guitar by aligning published scores to
recordings~\cite{riley2024domain}, then trained it on the GAPS classical guitar
dataset~\cite{riley2024gaps}, with state-of-the-art results on GuitarSet both
when trained on it and zero-shot \reported.

\textbf{Basic Pitch}~\cite{bittner2022basicpitch} is a small model that works
on any instrument. It outputs onsets, notes and a \term{pitch contour} at a
third of a semitone. It is fast and a useful baseline, and its contour is the
basis of our technique detection, but being instrument-agnostic it tends to
report string harmonics as extra notes.

\subsection{Electric guitar and tone}

Amplifier tone is the main obstacle. EGDB showed that accuracy depends strongly
on the amplifier~\cite{chen2022egdb}. EGDB-PG trained a Tone-informed
Transformer (TIT) on many amplifier-cabinet renders with a tone
embedding~\cite{chen2025egdbpg}; Pedroza et al.\ trained on synthetic
performances played through recorded tones and effects~\cite{pedroza2024tones}.
The recipe is the same in all of them: show the model the same performance in
many tones.

\subsection{Sequence-to-sequence transcribers}
\label{sec:notes-seq2seq}

Hawthorne et al.~\cite{hawthorne2021seq2seq} showed that a generic transformer
that ``translates'' a spectrogram into MIDI-like tokens matches specialised piano
models. MT3~\cite{gardner2022mt3} extended this to many instruments with tokens
that include the instrument. YourMT3+~\cite{chang2024yourmt3} improved the
encoder and the training mixtures but remains limited on pop recordings.
MuScriptor~\cite{rouard2026muscriptor} changes the data rather than the
architecture: a decoder-only transformer (up to 1.3\,B parameters, 5\,s
inputs) pre-trained on 1.45\,M synthesised MIDI files, trained on 170\,000
aligned real recordings, then refined with reinforcement learning. Distorted
electric guitar is the second most frequent class in its real data (about 47\%
of tracks). It can be told which instruments are present. Two limitations
matter for us: its tokens cannot represent the same pitch twice at once on one
instrument (keeping such notes in its test set lowers the large model's onset
F1 from 60.4 to 51.8 \reported), and like all models of this kind it outputs no
string and merges all guitars of one class. Its weights are non-commercial.

\begin{inapp}[: note models \done]
\begin{itemize}
\item \textbf{Guitar, default}: the GAPS CRNN (\file{transcribe/guitar_crnn.rs},
  optional download by \code{scripts/install_transcription_models.sh}, exported
  to ONNX by \code{scripts/export_guitar_crnn.py}, MIT). Audio is resampled to
  16\kHz{} and cut into 10\,s windows with 50\% overlap, run in parallel on at
  most 4 threads (each window takes about 600\,MB in tract). Notes are read
  with Kong's post-processing: onset and offset peaks with sub-frame shift, a
  frame threshold, notes under 30\ms{} dropped.
\item \textbf{Guitar, electric option}: \code{guitar_electric_crnn.onnx}, our
  fine-tune (``v2'', below), chosen in the Transcribe form or with
  \cmd{--guitar-model electric}. Local file only, because EGDB has no licence.
\item \textbf{Fallback and bass}: Basic Pitch (\file{transcribe/basic_pitch.rs},
  bundled, Apache-2.0; 22.05\kHz, 2\,s windows). It always transcribes the bass,
  and the guitar when the CRNN is not installed. With the CRNN it still runs on
  the guitar audio, only for its pitch contour (Section~\ref{sec:techniques}).
\item \code{prune_weak} removes weak simultaneous detections and, for Basic
  Pitch, likely overtones.
\end{itemize}
\end{inapp}

\begin{inapp}[: how v2 was trained]
\code{scripts/train/finetune_crnn.py} starts from the GAPS weights and trains on
EGDB (training ids), Guitar-TECHS chords and scales and GuitarSet players
00--02. Each DI recording is rendered through the app's own factory presets with
the desktop tool \code{render_tones} (\file{backend/offline.rs}); ``Metal High
Gain'' and ``Blues Crunch'' are left out of training because they are the
distortion benchmark. Labels are first shifted per recording to remove timing
offsets (\code{align_labels.py}). BatchNorm layers are frozen (batches are
small), precision is fp32/TF32 (fp16 overflowed in the power spectrum; bf16
crashed with NaNs at step 1100). The adopted run (\code{runs/spark_tf32}) took
8000 steps of 8 clips, about 1\,h on one DGX Spark; the best validation
checkpoint was step 5000.
\end{inapp}

\subsection{What we measured}

\begin{table}[h]
\centering
\caption{Note F50 (mean per item) on held-out sets \measured{}
(\code{bench_notes.py}, up to 12 items per group). MS = MuScriptor, conditioned
on the guitar classes; multi-instrument models are scored on their guitar
notes. Speed in audio seconds per second on one DGX Spark GPU (Basic Pitch on
CPU).}
\label{tab:bench-notes}
\footnotesize
\setlength{\tabcolsep}{4pt}
\begin{tabular}{@{}lccccccc@{}}
\toprule
Set & GAPS & v2 & Basic Pitch & YourMT3+ & MS small & MS medium & MS large \\
\midrule
EGDB, DI + 3 amp sims & 0.74 & \textbf{0.84} & 0.67 & 0.71 & 0.73 & 0.74 & 0.76 \\
EGDB real recordings (5) & 0.89 & \textbf{0.90} & 0.69 & 0.83 & 0.71 & 0.79 & 0.59 \\
EGSet12, real, miked amp & \textbf{0.89} & 0.88 & 0.76 & 0.54 & 0.85 & 0.88 & \textbf{0.89} \\
GOAT DI & 0.82 & \textbf{0.84} & 0.74 & 0.81 & 0.74 & 0.83 & 0.82 \\
GOAT, 5 amp renders & 0.56 & 0.71 & 0.56 & 0.39 & 0.66 & 0.76 & \textbf{0.77} \\
Guitar-TECHS excerpts & \textbf{0.81} & 0.78 & 0.61 & 0.70 & 0.66 & 0.73 & 0.71 \\
GuitarSet comp & 0.87 & 0.89 & 0.75 & \textbf{0.92} & 0.77 & 0.77 & 0.77 \\
GuitarSet solo & 0.91 & 0.92 & 0.81 & \textbf{0.93} & 0.86 & 0.89 & 0.86 \\
\midrule
Speed ($\times$ real time) & 34 & 89 & 165 & 14 & 6.4 & 4.0 & 1.3 \\
\bottomrule
\end{tabular}
\end{table}

How to read Table~\ref{tab:bench-notes}:
\begin{itemize}
\item \textbf{No model wins everywhere.} On amplified guitar the models trained
  on many tones lead: MuScriptor large 0.77 and v2 0.71 on the GOAT amplifier
  renders, against 0.56 for GAPS. MuScriptor large also makes the fewest
  octave-low errors there (2.1\% of notes against 8.0\% for GAPS).
\item \textbf{Real electric guitar through a clean amp} (EGSet12): GAPS, v2 and
  MuScriptor large tie at 0.88--0.89.
\item \textbf{YourMT3+} is best on acoustic GuitarSet and has the best note
  ends, but labels much amplified guitar as other instruments (0.39 on GOAT
  amps).
\item \textbf{Speed}: MuScriptor large runs at 1.3$\times$ real time on a large
  GPU, too slow for the app's CPU path; v2 runs at 89$\times$.
\end{itemize}

In the full app pipeline (\code{eval_guitarset}, GuitarSet players 03--05 played
through the app's presets, final-tab F1 median / octave errors) \measured:
\begin{center}
\small
\begin{tabular}{@{}lcc@{}}
\toprule
Tone & GAPS & v2 (electric) \\
\midrule
Clean & 0.84 / 1.2\% & \textbf{0.86} / 0.8\% \\
``Blues Crunch'' & 0.76 / 1.5\% & \textbf{0.81} / 1.0\% \\
``Metal High Gain'' & 0.43 / 4.5\% & \textbf{0.69} / 2.3\% \\
\bottomrule
\end{tabular}
\end{center}
A hand-written rule that raised notes struck with a strong partial a twelfth
above them (the distortion difference tone) made clean and crunch worse and
metal no better, so it was dropped: distortion needs a model trained on
distorted guitar, which is what v2 is.

\subsection{Failure modes}
\begin{itemize}
\item \textbf{Distortion}: octave-low errors from difference tones
  (Section~\ref{sec:audio-harmonics}). The 8-string SoundFont render of
  ``Meshuggah -- Bleed'' gets a final-tab F1 of 0.02 with the GAPS CRNN, against
  0.24 with Basic Pitch \measured.
\item \textbf{Distorted guitar in real mixes}: on MoisesDB, v2 on the
  BS-RoFormer stem scores 0.56 on distorted guitars, 0.75 on clean electric and
  0.83 on acoustic \measured. This is the largest remaining loss on real songs.
\item \textbf{Separation artefacts} are unseen in training (no published guitar
  transcriber is trained on separated stems); planned as training task T1
  (Section~\ref{sec:next}).
\item \textbf{Same-pitch notes} (unisons on two strings) cannot be represented
  by MT3-style tokens.
\end{itemize}
